\section*{Generative architecture and conservation of information}
\addcontentsline{toc}{section}{Generative architecture and conservation of information}
\label{libnar:troncal}

Triadic Modular Holography, abbreviated HMT, locates the origin of the
construction in a finite system of positions, operations, and transports.
Its prolongation yields compatible histories; the histories yield
numerical and incidence readings; these realizations yield the structures
with which Dimensional Mechanics operates. The question organizing this
article concerns that constructive continuity: what remains when a
unit changes scale, when an operation modifies its representation,
and when part of the state is transferred to memory.

The content of the unit is specified throughout this progression. In an
arithmetic partition, the normalized sum of its components is conserved.
In an inner-product space, norms and correlations are conserved.
In a geometric representation, forms, directions, and incidences
are transported. Each conservation has its own object and equation.
Their articulation makes it possible to follow a single construction from the carry
of a digit to the recovery of a state distributed across an archive
of arbitrary depth.

The following exposition presents this progression in dependency order.
The complete definitions and proofs appear later, in the
mathematical body. The purpose here is to enable the reader to recognize
the role of each operation before studying its general formulation: how a
history is produced, what determines a reading, and how information
that a visible representation no longer displays can be reconstructed.

\subsection*{From finite positions to operational histories}
\label{libnar:app}

The construction called APP, for \emph{All Possible Paths}, begins with
two axes of nine marks and two operations on their pairs of
positions: addition and multiplication. The positions have an identity that is
preserved when they are evaluated. The pair \((i,j)\) admits the additive reading
\(i+j\) and the multiplicative reading \(ij\); a route also records how that
pair was reached, which operation acted, and which orientation the
transport followed. The relation between additive accumulation and
multiplicative composition is thus situated on a common support.

The modular step preserves both the remainder and the quotient. For a
positive evaluation \(n\), the nonadic mark \(r\), lying between
one and nine, and the quotient \(q\) satisfy
\[
 n=r+9q.
\]
This equality allows the integer evaluation to be recovered from its enriched
modular reading. Two positions may share the same mark while
retaining different quotients. For example, \((5,5)\) and \((8,2)\)
produce the sum ten; their products are twenty-five and sixteen. Both
products have mark seven, with quotients two and one, respectively.
The operational difference is located in the record accompanying
the mark. The definitions in \cref{sec:nucleo} develop this
conservation jointly with the operation sheets and routes.

The oriented ternary unit, called TRIT, distinguishes three regimes
by means of the signs \(-1,0,+1\). Its action includes phase selection,
the orientation of prolongations, and carry. Balanced ternary
division writes an integer as \(n=r+3c\), with
\(r\in\{-1,0,+1\}\); the remainder identifies the local state and the
quotient preserves what is transported to the next level. Subsequent
geometric realizations distinguish the elliptic, parabolic, and
hyperbolic regimes. In the temporal reading of the corpus, return with memory,
the threshold, and bidirectional opening have specific actions on
the state.

The \emph{Topological Phase Kernel}, or topological phase kernel, is
abbreviated TPK. It composes tritic selection, transport, and
memory updating. A transition modifies positions, sheets, or
phase and preserves the data needed to relate the previous state
to the next. Its result is an enriched state: it contains
residues, quotients, carries, orientation, boundary, and route prefixes.
The number that will subsequently be read arises from that operational
history. Its genealogy preserves the operations that allow it
to be continued and the relations that allow it to be compared with other
histories.

This organization gives return a precise meaning. The calendar
of one hundred and eight positions is distributed across twelve windows of nine.
Completing a nonadic turn returns the corresponding phase and advances
the window record. Returning to a visible position and
recovering the enriched state are therefore two distinct events.
\cref{sec:extension} formalizes the calendar,
holonomy, and the memory increment accompanying their returns.

\subsection*{Contraction of the value and persistence of incidence}
\label{libnar:continuo}

A history is prolonged through compatible prefixes. Each prefix
preserves the operations performed and delimits its admissible
continuations. The joint construction of the continuum maintains graded
transport, memory balances, ternary incidences, boundary relations,
and determinantal compositions within that same system.
These operations are linked through their actions on
the same histories and through their compatibility with refinement
(\cref{iv:continuo-conjunto}). Their correlation is part of
the construction of the object.

The Archimedean reading associates cylinders of increasing precision
with the prefixes. Their compatible intervals contract and determine a value
when their diameters tend to zero. The half-open convention distinguishes
boundary representatives during refinement; carry
information records the continuations that reach the same endpoint.
This reading gives the numerical realization of the generated continuum.
The sequence of intervals expresses how much of the value has been determined;
the enriched state expresses how it has been determined.

The incidence reading preserves other information: which elements form
a block, how a boundary is oriented, which transports relate the
prefixes, and which actions preserve those relations. Through the maps
developed in \cref{exc:seccion}, this information is realized in
Paley and Hadamard matrices, Witt designs and Steiner systems,
Golay codes, lattice constructions, and the Leech lattice.
Prolongation to the vertex operator algebra \(V^\natural\)
locates the action of the Monster group there. Each step preserves its carrier,
its map, and its compatibility conditions.

The relation between number and symmetry thus becomes concrete. The
Archimedean value identifies histories under a numerical equivalence; the
incidence structure preserves relations between histories and their boundaries.
Both readings arise from the same state and fulfil complementary
functions. A scalar magnitude may remain fixed while
the incidence representation undergoes a transformation. Following
both readings makes it possible to study jointly the stabilization of
the value and the evolution of the structure realizing it.

The unit is also conserved when distributed among prolongations.
The measures of the successor cylinders sum to the measure of the predecessor cylinder.
In the associated function spaces, this equality becomes
a refinement isometry: a function from the previous level preserves
its norm when represented on the new positions. The same
principle relates partition, measure, and transport. The proofs
of the joint construction specify composition and passage to the
limit, so that every new depth remains linked to the
preceding ones.

The realization of a measurement adds a concrete operation to this relation
between state and reading. The system and the device couple, produce
a record, and establish which distinctions become accessible. In an
isometric realization of the coupling, the joint description preserves
the correlations; the recorded outcome then determines the
corresponding update. Evaluating an observable, conditioning the
state on an outcome, and preparing the record anew have distinct
functions. This sequence makes it possible to follow both the information that
reaches the reading and that which remains in the state and the device.
Memory intervenes in this operational continuity between an observation
and subsequent ones (Sections~\ref{nuclear:3} and~\ref{lib04:archivo}).

\subsection*{The dimensional unit and recoverable carry}
\label{libnar:unidad}

The holographic extreme and mean ratio is studied here through the
operations that conserve the unit during its extension and
refinement. The added position has its own identity. If
\(E_9\) contains the nine nonadic classes, its completion is written
\(E_9\sqcup\{\ast\}\). The zero residue class already belongs to
\(E_9\) and has the positive representative nine; \(\ast\) denotes
the additional boundary state. The tenth possibility thus retains a
function distinct from modular reduction. In several coordinates, counting
how many have reached \(\ast\) determines the strata of the completion
and makes it possible to follow each stratum's contribution to the unit.

The cubic completion provides a first exact
representation. Extending the positions of each coordinate from nine to ten
decomposes the cubic unit into the interior, extension strata, and
final vertex:
\[
 10^3=9^3+3\cdot9^2+3\cdot9+1
     =729+243+27+1=729+271.
\]
Dividing by one thousand converts the count into a partition of unit mass.
The terms preserve their dimensional provenance: one, two, or three
coordinates occupy the added position. The expansion
\((9+1)^d\) preserves this distinction for every positive
integer dimension. Increasing dimension and changing scale are
related by coordinate maps and sums of strata
(\cref{sec:revision-emrh,lib03:completacion}).

Carry introduces a finer transformation. Consider the
partition \(73+27=100\). Multiplication by ten produces
\(730+270=1000\). Transferring one unit from the first component to
the second gives
\[
 (73,27,1)\longmapsto(10\cdot73-1,10\cdot27+1)
                  =(729,271).
\]
The third datum is the carry. The final sum remains one thousand and the
normalized partition remains one. At the same time, the decimal residue
of the second component preserves the transfer: from \(271\) one
recovers the remainder one, then \((271-1)/10=27\), and finally
\((729+1)/10=73\). The result carries with it the data of the
operation that produced it.

Theorem~\ref{lib03:teo-residuo} proves this inversion for the map
\[
 L_c(x,y)=(Bx-c,By+c),\qquad x+y=M,
\]
with base \(B\), canonical carry \(0\leq c<B\), and the stated
positivity conditions. After several stages, the accumulated record
\(C_n\) collects the carries as a word in base \(B\), and the
components satisfy
\[
 x_n=B^n x_0-C_n,\qquad y_n=B^n y_0+C_n.
\]
The sum preserves \(B^nM\). With base and depth fixed, the final
state recovers the complete carry word and the initial partition
(\cref{lib03:cor-palabra}). Depth is part of the reconstruction
data, since it also determines the initial positions that
contain zeros.

The normalized readings converge with explicit tail bounds.
Three levels of description are thereby linked: a reversible integer
operation, a sequence of partitions of the unit, and a system of
cylinders determining a limit reading. The effective selection of
continuations follows the rules of the transport under consideration. The example
\((729,271)\) shows its first transfer; the refinement
proofs also describe the memory of all subsequent transfers
and the boundary coincidences between expansions
(\cref{lib03:teo-limite,lib03:teo-cilindros}).

\subsection*{Numerical regions and the coupling register}
\label{libnar:registro}

In the corpus, the coordinates \(\pi_{\rm HMT}\),
\(\varphi_{\rm HMT}\), and \(e_{\rm HMT}\) arise from regional
readers on the APP--TRIT--TPK prolongation. The closure region,
propagation, and stable relation between modes have their own
operators. For example, the calendar of additive and multiplicative modes
produces the incidence matrix
\(F_{\rm av}=\bigl(\begin{smallmatrix}0&1\\1&1\end{smallmatrix}\bigr)\).
Its stable positive direction determines the golden coordinate in the
corresponding realization. The generation proofs specify, for
each requested depth, how a certified prefix is obtained and
how its compatibility with subsequent prefixes is preserved
(\cref{sec:generacion}).

The same terminal state preserves an oriented register of its twelve
windows. Extraction of the signed channels, total charge, and
Hadamard inversion produce the ordered vector \(K\). The term
\emph{holographic arithmetic--geometric coupling constant}
denotes that register together with its readers and its fixed phase origin.
Its function is to relate an exact arithmetic encoding
to the incidence and geometric realizations acting on the
same object (\cref{lib01:definicion}).

In the canonical integral representation one obtains
\[
 \begin{aligned}
 K=\bigl(&234,543,140,729,659,824,\\
         &621,58,914,794,146,601\bigr),
 \qquad\sum_{j=1}^{12}K_j=6263.
 \end{aligned}
\]
\cref{lib01:registro} shows the complete extraction, the
congruences of its channels, and the divisibility allowing the Hadamard
matrix to be inverted. The order of its components preserves
window position. This information intervenes in the digit
reading as well as in cyclic operators and incidences.

The total sum serves a reconstruction function. Differences
between positions separated by three and four steps jointly
determine the relative variations of the register. Adding the same
quantity to all twelve components preserves these differences; the sum fixes
precisely that common level. In \(K\), the charge \(Q(K)=6263\)
restores the uniform component. Differences, charge, and integer
compatibilities allow the complete register to be recovered
(\cref{prop:k-transversal}). The global relation between positions and
the relative organization of their channels are preserved by
complementary readings.

In base one thousand, with twelve declared positions, the register is encoded
by the repeating fraction
\[
 \kappa_{\rm ag}=
 \frac{234543140729659824621058914794146601}{10^{36}-1}.
\]
Each component occupies three digits; thus, the component fifty-eight
is written \(058\) in the expansion. Multiplying the fraction
by its denominator and performing twelve Euclidean divisions recovers the
exact register. The periodicity of this encoding describes a
fixed reading of the finite register; updating the enriched
history, for its part, preserves its own memory dynamics.

The connection with the coordinate \(\alpha_A\) is expressed in the
closure reader by
\[
 \alpha_A=
 \pi_{\rm HMT}+e_{\rm HMT}-\varphi_{\rm HMT}-4-\kappa_{\rm ag}.
\]
The three regional coordinates and the register arise from the generated
state and act in the constructive order described in
\cref{lib01:clausura}. The formula makes their arithmetic coupling
visible: the extraction of \(K\) preserves part of the
structure intervening in that closure. The other realizations
of \(K\) preserve their respective readers. In this way, a
scalar identity, an incidence matrix, and a phase transport
can be studied from the same provenance, with their actions
clearly distinguished.

\subsection*{A register capable of identifying a transformation}
\label{libnar:identificacion}

The function of the register can be examined through a concrete
operation. Let \(S\) be the cyclic shift of its twelve
components. The twelve columns
\(K,SK,\ldots,S^{11}K\) form an invertible matrix \(O_K\).
The result means that the cyclic orbit of this register reaches
all directions in its twelve-component space. The proof
computes the spectrum of the Gram matrix \(O_K^*O_K\) and gives exact
positive bounds for its eigenvalues (\cref{lib01:marco-ciclico}).

Suppose that an operator \(H\) respects the cyclic shift,
that is, \(HS=SH\). When the vector response \(y=HK\) is observed,
its shifts provide the responses to all columns
of \(O_K\). The complete operator is then recovered by
\[
 H=O_yO_K^{-1},\qquad
 O_y=(y,Sy,\ldots,S^{11}y).
\]
The observation contains twelve components, and the commutation
hypothesis identifies the class of recoverable transformations.
Within that class, a single vector response determines all
coefficients of the circular filter. For example, the combination
\(2I+3S-S^5\) is identified from its action on \(K\),
with the error bounds of Theorem~\ref{lib01:operadores}.

This fact provides an operational application of the structure of the
register. Its order and spectrum allow it to be used as a test
state for identifying compatible transformations. If the state
and its responses are transported by a full-memory isometry,
the Gram matrix of the orbit is conserved and identification remains
recoverable by the adjoint. Stability follows from computed
eigenvalues and explicit inverse operators. The same
register participating in arithmetic closure thus acquires a
precise function in the recognition of transports.

Normalization also has associated information. The
unitary factor of \(O_K\) preserves orientation and angles; its Gram matrix
preserves the amplitudes needed to reconstruct the original
matrix. The pair formed by the two factors recovers \(O_K\)
(\cref{lib01:polar}). This separation allows a normalized
representation to be chosen while maintaining the exact
scale of the register.

\subsection*{Two transports and bilateral conservation}
\label{libnar:bilateral}

The next step relates refinement to a law of
correlation conservation. Let \(\mathcal B\) and
\(\mathcal C\) be two isometries transporting the same state space
to a common output space. For \(a>0\), form the
two operators
\[
 Q_-=a\mathcal B-\mathcal C,\qquad
 Q_+=(a+1)\mathcal B+\mathcal C.
\]
The two combinations bring together a common contribution and a
relative contribution with different signs. Summing their norms
with the corresponding weights cancels the cross terms:
\[
 (a+1)\|Q_-v\|^2+a\|Q_+v\|^2
   =(2a+1)(a^2+a+1)\|v\|^2.
\]
Theorem~\ref{lib04:teo-balance} proves this equality for spaces of any
dimension and for the stated isometric
transports. The same cancellation preserves inner products and
admits a formulation for forms preserved by the transports.

In the case \(a=8\), the factor \(a^2+a+1\) equals seventy-three.
The factorization of the full isometry first distributes the
norm between two components with weights
\[
 \frac{72}{73}\qquad\text{and}\qquad\frac{1}{73}.
\]
An orthogonal mixing then acts to produce the normalized
channels of \(Q_-\) and \(Q_+\). The weights thus have a
precise operational origin: \(72=8\cdot9\) and
\(73=8\cdot9+1\). Their function is preserved when the dimension
or the transports change, within the hypotheses of the bilateral law
(\cref{lib04:factorizacion}).

Recovery is seen directly in the output combinations.
If \(x=Q_-v\) and \(y=Q_+v\), then
\[
 x+y=(2a+1)\mathcal Bv,\qquad
 ay-(a+1)x=(2a+1)\mathcal Cv.
\]
The two transports of the initial state remain determined by
the pair of channels. The adjoint of either isometry
then recovers the corresponding state. With full
normalization, the operator \(W_a\) satisfies \(W_a^*W_a=I\):
it preserves the norm and provides an exact inverse on its image.

A particular reading may vary while all this
information is conserved. When \(\mathcal C=\mathcal B R\), with
\(\mathcal B\) and \(R\) unitary, the latter serving as relative transport, the
diagonal reading of an observable \(G\) in both channels produces
\[
 \mathcal T(G)=\frac{72G+R^*GR}{73}
 \qquad(a=8).
\]
Observables compatible with \(R\) remain fixed; the others
record the difference between its two actions. This equation distinguishes
conservation of the complete state from the response of a particular
reader (\cref{lib04:teo-lector}).

\subsection*{The same carry, now as a transportable state}
\label{libnar:ejemplo}

The preceding relation applies materially to carry. Each
admissible label \((x,y,c)\) is associated with a vector of an
orthonormal basis. The integer transformation \(L_c\) then induces the
operator \(J\) taking that basis to the output labels
\((Bx-c,By+c)\). Its inverse recovers the carry through the
residue. The bijection of labels makes \(J\) a
unitary transformation (\cref{lib05:prop-unitaria}). Integer
values designate the labels; norms correspond to the
amplitudes of the vectors representing them.

For base ten and initial sum one hundred, consider the two labels
\(d_1=(73,27,1)\) and \(d_2=(73,27,2)\). Their images are
\(e_1=(729,271)\) and \(e_2=(728,272)\). Fix the transport
\(R\) that interchanges these two labels and leaves the others fixed.
The composition uses \(\mathcal B=J\) and \(\mathcal C=JR\).
This choice provides an explicit example within the bilateral
theorem, with the transport rule specified before its evaluation.

On the state \(|d_1\rangle\), the channels for \(a=8\)
are
\[
 Q_-|d_1\rangle=8|e_1\rangle-|e_2\rangle,
 \qquad
 Q_+|d_1\rangle=9|e_1\rangle+|e_2\rangle.
\]
Their squared norms are sixty-five and eighty-two. The exact
balance is \(9\cdot65+8\cdot82=1241=17\cdot73\), while
the sum of the two channels recovers \(17|e_1\rangle\).
Applying the label inverse recovers the partition \((73,27)\)
and carry one.

The fractional reading of the first component also provides an
exact prediction for this declared transport. Evaluating it
diagonally in both normalized channels yields
\[
 \frac{72\cdot729+728}{73\cdot1000}
   =\frac{729}{1000}-\frac{1}{73000}.
\]
The diagonal carry reader gives \((72\cdot1+2)/73=74/73\).
The observable transported from the input, by contrast, preserves
the eigenvalue one of the original carry. The difference between the two
results makes it possible to identify which reading operation has
been performed. The state information remains recoverable, and the
diagonal reading quantifies how the relative transport intervenes
(\cref{lib05:ejemplo}).

The example brings together the two conservations of the article. The
arithmetic sum is conserved in the refined labels, and the norm is
conserved in their unitary representations. The connection between them
arises from the bijection and its lift to orthonormal bases.
This composition allows readers to be computed, transport to be inverted,
and error to be controlled while keeping the objects on
which each law acts identified.

\subsection*{Full memory, symmetries, and arbitrary depth}
\label{libnar:memoria}

Bilateral conservation is prolonged stage by stage. At each level,
one channel continues the state and the other records a memory
complement. If \(R_Nv\) is the state reaching level \(N\) and
\(m_0,\ldots,m_{N-1}\) are the recorded complements, the
accumulated balance is
\[
 \|v\|^2=\|R_Nv\|^2+\sum_{n=0}^{N-1}\|m_n\|^2.
\]
The equality identifies exactly where the initial norm
remains. Finite recovery uses the continued state and
all available complements; their correlations are preserved
by the same joint isometry.

For the archiving policy of Theorem~\ref{lib04:teo-archivo}, the
norm of the continued channel has a uniform geometric bound.
In the case \(a=8\), its squared norm is bounded at every stage
by \(729/1241\). As depth increases, the continued
component tends to zero and the infinite archive preserves the entire
norm. If \(A_n\) and \(D_n\) denote the continuation
and memory blocks, respectively, the state is reconstructed
by the convergent series
\[
 v=\sum_{n\geq0}R_n^*D_n^*m_n.
\]
The inverse operator has norm one. A perturbation of norm
\(\varepsilon\) in the complete archive produces a reconstruction
error of norm at most \(\varepsilon\). Finite
truncations have their own bounds and, when only
their memory is used, the finite inverse specified in the same
development.

\pagebreak % REV03: balance the complete concluding narrative paragraphs.
The incidence structure of the register \(K\), its cyclic orbit, and its
distinguished directions can be transported through these
isometries. The inner products determining angles
and Gram matrices are preserved; projectors and readers are
transported with their domains; the action of the adjoint recovers the
initial representation. Exterior powers transport
orientation and preserve Gram determinants, which measure the
squares of the corresponding volumes. The dimension of the
carrier may be increased, and the geometric information of the
image remains specified by the isometry.

The relation to Noether conservation is formulated at a
well-defined additional level. For the auxiliary discrete action
constructed in \cref{lib01:carga-noether}, a collection of
generators intertwining the transports produces a
conserved charge of the form \(p_n^*A_nq_n\). A variational
action, its equations, and the symmetry preserving them intervene here.
The norm balance and the variational charge are articulated through
the common transports, each with its hypotheses and its
function in the proof.

This progression establishes an operational relation between unit,
number, memory, and geometry. The arithmetic unit admits
reversible refinements; the enriched state preserves the
record of operations; the numerical and incidence
realizations extract correlated properties; bilateral coupling
maintains information and correlations during their
redistribution. The holographic extreme and mean ratio thus acquires
a demonstrative development in which partition, extension,
transport, and recovery can be followed through explicit
maps. The following pages present those maps,
their proofs, and their complete examples.
